% ~1 página
\section{Introdução}
\label{sec:introducao}

Os sistemas de condução autónoma organizam-se, em geral, numa cadeia de
módulos: perceção, previsão do movimento dos outros utentes da via,
planeamento e controlo~\cite{paden2016survey,schwarting2018planning}. A
previsão e o planeamento dependem um do outro. Um plano só é seguro se
contar com o sítio onde os outros veículos vão estar nos segundos
seguintes, e esse futuro não é conhecido: tem de ser estimado a partir do
que se observou até ao momento.

Este trabalho junta as duas metades do problema. Na previsão, a questão é
quanto se ganha com modelos de aprendizagem automática clássica (árvores
de decisão, florestas aleatórias, \emph{gradient boosting}) em relação às
extrapolações físicas mais simples. A restrição a modelos clássicos é
deliberada: treinam-se em segundos num computador comum, funcionam com
poucos dados e deixam ver que variáveis pesam na previsão. No
planeamento, a questão é se uma previsão melhor dá um plano melhor,
quando as previsões entram como obstáculos móveis num planeador de
velocidade baseado em pesquisa A*.

Os dados são os do Waymo Open Motion Dataset
(WOMD)~\cite{ettinger2021large}, que regista trajetórias reais de
veículos, peões e ciclistas em meio urbano. A escolha do problema foi
condicionada pelo que os dados permitem. Cada cena dura cerca de
\un{20}{s}, o veículo da Waymo percorre nela \un{96}{m} (mediana) e, em
82\% das cenas analisadas, há um único caminho possível entre o ponto de
partida e o de chegada. Não há pedidos de viagem nem frotas. Problemas de
rotas longas ou de despacho de veículos exigiriam dados inventados; as
interações locais entre veículos, pelo contrário, estão presentes em
todas as cenas.

O artigo tem três contributos:
\begin{itemize}
  \item uma revisão do estado da arte em previsão de trajetórias, em
    planeamento de movimento e na ligação entre os dois
    (\cref{sec:estado-arte});
  \item uma arquitetura «prever para planear» com modelos clássicos e um
    planeador A* (\cref{sec:abordagem});
  \item resultados preliminares em 149 cenas reais, com as limitações que
    ainda têm (\cref{sec:resultados}).
\end{itemize}
